% GENERATED FILE. Edit canonical lessons, then run npm run generate:books.
% canonical-source-hash: fnv1a64:0fc45ac8ebcaa73f
% canonical-lessons: ZH-C278-zuihou, ZH-W278-zui, ZH-C278-sequence, ZH-W278-ran, ZH-C278-shouxian, ZH-C278-jiezhe, ZH-C278-le-yihou, ZH-C278-practice

\chapter{First, Then, Last}
\label{ch:zh-first-then-last}
\hlchaptermodality{278}

\begin{chapteropening}
\textbf{By the end of this chapter:} \emph{I can tell what happened first, next and last.}

Telling things in order with \zh{先}, \zh{然后} and \zh{最后}, first of all and straight after, and \zh{了} … \zh{以后} for after doing, built on two new written characters: \zh{最} and \zh{然}.
\end{chapteropening}

\section[\texorpdfstring{\zh{最后} (zuìhòu)}{最后 (zuìhòu)}]{\zh{最后} (zuìhòu) — last; finally}

\label{lesson:ZH-C278-zuihou}



\begin{quote}
A new chapter, and it opens on the word for the end of a story.

\begin{itemize}
\raggedright
  \item \emph{Recall:} say \emph{rè'ài}, to love deeply — \textbf{R1}, the chapter you have finished
  \item \emph{Recall:} say \emph{xiǎngshòu}, to enjoy — \textbf{R1}
  \item \emph{Recall:} say \emph{zuì}, most — \textbf{R3}, durable
  \item \emph{Recall:} say \emph{hòu}, after; behind — \textbf{R3}
\end{itemize}
\end{quote}

\subsection*{You'll want to know: \zh{最后}}
\textbf{\zh{最后}} — \emph{zuìhòu} — \textquotedblleft{}last; finally; in the end\textquotedblright{}.

\begin{quote}
\textbf{\zh{最后我们回家了}} \emph{zuìhòu wǒmen huí jiā le} — In the end we went home.
\end{quote}

\begin{quote}
\textbf{\zh{今天是最后一天}} \emph{jīntiān shì zuìhòu yì tiān} — Today is the last day.
\end{quote}

The next lesson puts the first character in your hand. It is the \emph{zuì} you already say for most.

\begin{cousinweb}
\textbf{\zh{最后}} (\emph{zuìhòu}) is built from \zh{最} 'most' and \zh{后} 'after', both inherited from Middle Chinese, \emph{tswajH} and \emph{huwX} (Guangyun, in Baxter's transcription). The dictionary glosses the first 'utmost' and lists the second, under its traditional form, with 'before and after'; 'last' for the pair is the word's own sense.
\end{cousinweb}

\begin{grammarlens}[title={What you've built}]
\zh{最后} opens the last step of a story: \textbf{\zh{最后我们回家了}} \emph{zuìhòu wǒmen huí jiā le}. Before a noun it picks out the last one: \textbf{\zh{最后一天}} \emph{zuìhòu yì tiān}.

The \zh{了} at the end is the \zh{了} you know: it happened.
\end{grammarlens}

\subsection*{Your turn}
\begin{itemize}
\raggedright
  \item \emph{Say it:} \emph{zuìhòu}
  \item \emph{Say it:} \emph{zuìhòu}, once more
  \item \emph{Say it:} \zh{最后我们回家了} — \emph{zuìhòu wǒmen huí jiā le} — In the end we went home.
  \item \emph{Say it:} \zh{今天是最后一天} — \emph{jīntiān shì zuìhòu yì tiān} — Today is the last day.
  \item \emph{Recall:} say \emph{zuì}, then \emph{zuìhòu}
\end{itemize}

\begin{tcolorbox}[breakable,colback=teal!4,colframe=teal!35!black,title={Before you move on}]
What does \zh{最后} mean? (\textquotedblleft{}Last; finally\textquotedblright{}.) Which word you know is inside it? (\emph{zuì}, most.)
\end{tcolorbox}
\section[\texorpdfstring{\zh{最} (zuì)}{最 (zuì)}]{\zh{最} — \zh{日}, six strokes, then \zh{又}}

\label{lesson:ZH-W278-zui}



\begin{quote}
The word first, then the character inside it.

\begin{itemize}
\raggedright
  \item \emph{Recall:} say \emph{zuìhòu} — \textbf{R1}, one lesson back, before it decays
\end{itemize}

\begin{quote}
\zh{日} \zh{对}
\end{quote}

\emph{Read it:} each one aloud, with nothing said first — \emph{rì}, \emph{duì}
\end{quote}

\subsection*{Script — \zh{日}, six strokes, then \zh{又}}
\hlblockfigure{\detokenize{figures/ZH-W278-zui-filmstrip.pdf}}{How \zh{最} is written, stroke by stroke}

\begin{quote}
\zh{最}
\end{quote}

\textbf{Three pieces, and you write two of them already.}

\textbf{\zh{日}, four strokes}, wide and flat on top: the left side, the top that turns down the right side, the middle bar, the bottom bar.

\textbf{Six strokes under it, on the left}: a long horizontal across the whole width; two verticals down from it, the left one first; two short horizontals between them, the upper one first; and last a short stroke at the foot that rises to the right.

\textbf{\zh{又}, two strokes}, on the right, as on the left of \zh{对}: a horizontal that turns and falls to the lower left, then a long right-falling stroke.

4 + 6 + 2 = twelve.

\subsection*{Your turn}
\begin{itemize}
\raggedright
  \item \emph{Write it:} \zh{最}, all twelve, with the model in front of you
  \item \emph{Cover:} write it again from memory, and count the pieces, not the strokes
\end{itemize}

\begin{tcolorbox}[breakable,colback=teal!4,colframe=teal!35!black,title={Before you move on}]
How many strokes in \zh{最}? (\textbf{Twelve}.) What sits on top? (\textbf{\zh{日}}.) Which stroke rises to the right? (\textbf{The last of the six, at the foot, before \zh{又}}.)

Source: \href{https://raw.githubusercontent.com/chanind/hanzi-writer-data/68d10a4b21150cae5e1ebbd223eed289cf32d90c/data/\%E6\%9C\%80.json}{Hanzi Writer Data \zh{最}, pinned stroke paths and medians}.
\end{tcolorbox}
\section[\texorpdfstring{\zh{先}…\zh{然后}…\zh{最后} (xiān … ránhòu … zuìhòu)}{先…然后…最后 (xiān … ránhòu … zuìhòu)}]{\zh{先}…\zh{然后}…\zh{最后} — first, then, last}

\label{lesson:ZH-C278-sequence}



\begin{quote}
Three words you have, and one story to put them in.

\begin{itemize}
\raggedright
  \item \emph{Read it:} \zh{最后}, aloud — \emph{zuìhòu}
  \item \emph{Recall:} say \zh{先} — \emph{xiān}, first — \textbf{R3}, durable
  \item \emph{Recall:} say \emph{ránhòu}, then; after that — \textbf{R3}
\end{itemize}
\end{quote}

\begin{grammarlens}[title={Grammar Lens: \zh{先}, \zh{然后}, \zh{最后} — first, then, last}]
\begin{quote}
\textbf{To tell things in order, put \zh{先} before the first verb, \zh{然后} at the front of the next step, and \zh{最后} at the front of the last.}
\end{quote}

\begin{quote}
\textbf{\zh{我先吃饭} \zh{然后看书} \zh{最后回家}} \emph{wǒ xiān chī fàn, ránhòu kàn shū, zuìhòu huí jiā} — First I eat, then I read, and last I go home.
\end{quote}

\zh{然后} is the \emph{ránhòu} you have said since the chapter on the past. Now it prints: the next lesson puts \zh{然} in your hand.

\zh{先} sits after the person, right before the verb: \zh{我先吃饭}. \zh{然后} and \zh{最后} open their clause, and the person need not be said again.

\begin{quote}
\textbf{\zh{他先去学校} \zh{然后去打工}} \emph{tā xiān qù xuéxiào, ránhòu qù dǎgōng} — He goes to school first, then to his part-time job.
\end{quote}

The order of the clauses is the order of the events.
\end{grammarlens}

\subsection*{Your turn}
Say these aloud:

\begin{itemize}
\raggedright
  \item \zh{我先吃饭} \zh{然后看书} \zh{最后回家} — \emph{wǒ xiān chī fàn, ránhòu kàn shū, zuìhòu huí jiā} — First I eat, then I read, and last I go home.
  \item \zh{他先去学校} \zh{然后去打工} — \emph{tā xiān qù xuéxiào, ránhòu qù dǎgōng} — He goes to school first, then to his part-time job.
  \item three things you did this morning, in order, with \zh{先}, \zh{然后} and \zh{最后}
\end{itemize}

\begin{tcolorbox}[breakable,colback=teal!4,colframe=teal!35!black,title={Before you move on}]
Where does \zh{先} go? (\textbf{After the person, before the first verb}.) And \zh{然后} and \zh{最后}? (\textbf{At the front of their clause}.)
\end{tcolorbox}
\section[\texorpdfstring{\zh{然} (rán)}{然 (rán)}]{\zh{然} — a four-stroke piece, \zh{大} with a dot, and four dots}

\label{lesson:ZH-W278-ran}



\begin{quote}
The pattern first, then the character it was waiting for.

\begin{itemize}
\raggedright
  \item \emph{Recall:} say \zh{我先吃饭} \zh{然后看书} — \emph{wǒ xiān chī fàn, ránhòu kàn shū} — \textbf{R1}, one lesson back, before it decays
\end{itemize}

\begin{quote}
\zh{大} \zh{点} \zh{热}
\end{quote}

\emph{Read it:} each one aloud, with nothing said first — \emph{dà}, \emph{diǎn}, \emph{rè}
\end{quote}

\subsection*{Script — a four-stroke piece, \zh{大} with a dot, and four dots}
\hlblockfigure{\detokenize{figures/ZH-W278-ran-filmstrip.pdf}}{How \zh{然} is written, stroke by stroke}

\begin{quote}
\zh{然}
\end{quote}

\textbf{Three pieces, and you write the last two already.}

\textbf{The upper left, four strokes}: a left-falling stroke; a horizontal that turns and falls to the lower left; then two dots inside it, the upper one first. It is not \zh{月}: there is no hooked frame and no inner bar.

\textbf{\zh{大} with a dot, four strokes}, on the upper right: the bar, the left-falling stroke from above it, the right-falling stroke, and last a dot at its upper right.

\textbf{Four dots along the foot}, from left to right, as at the foot of \zh{点} and \zh{热}.

4 + 4 + 4 = twelve.

\subsection*{Your turn}
\begin{itemize}
\raggedright
  \item \emph{Write it:} \zh{然}, all twelve, with the model in front of you
  \item \emph{Cover:} write it again from memory, and count the pieces, not the strokes
\end{itemize}

\begin{tcolorbox}[breakable,colback=teal!4,colframe=teal!35!black,title={Before you move on}]
How many strokes in \zh{然}? (\textbf{Twelve}.) Which character you write sits on the upper right? (\textbf{\zh{大}, with a dot added}.) Where else have you written the four dots? (\textbf{At the foot of \zh{点} and \zh{热}}.)

Source: \href{https://raw.githubusercontent.com/chanind/hanzi-writer-data/68d10a4b21150cae5e1ebbd223eed289cf32d90c/data/\%E7\%84\%B6.json}{Hanzi Writer Data \zh{然}, pinned stroke paths and medians}.
\end{tcolorbox}
\section[shǒuxiān]{shǒuxiān — first of all}

\label{lesson:ZH-C278-shouxian}



\begin{quote}
Before the new one: the character you wrote last, inside the word it belongs to.

\begin{itemize}
\raggedright
  \item \emph{Read it:} \zh{然后}, aloud — \emph{ránhòu}
  \item \emph{Recall:} say \zh{我先吃饭} — \emph{wǒ xiān chī fàn} — \textbf{R1}, two lessons back
\end{itemize}
\end{quote}

\subsection*{You'll want to know: shǒuxiān}
\emph{shǒuxiān} — \textquotedblleft{}first of all; to begin with\textquotedblright{}.

\begin{quote}
\emph{shǒuxiān wǒ xiǎng shuō xièxie} — First of all, I'd like to say thank you.
\end{quote}

\begin{quote}
\emph{shǒuxiān wǒmen qù mǎi dōngxi} — To begin with, we go shopping.
\end{quote}

Its first character comes later; for now it is a word for the ear.

\begin{cousinweb}
\emph{shǒuxiān} is built from shǒu 'the head' and \zh{先} 'first', both inherited from Middle Chinese, \emph{syuwX} and \emph{sen} (Guangyun, in Baxter's transcription). The dictionary glosses the first 'the head' and 'the beginning', and the second 'before and after'.
\end{cousinweb}

\begin{grammarlens}[title={What you've built}]
\zh{先} sits inside the clause, after the person: \zh{我先吃饭}. \emph{shǒuxiān} opens the sentence, before the person, and it sounds planned: the first of several points.

\zh{然后} and \zh{最后} follow it as usual: \emph{shǒuxiān …, ránhòu …, zuìhòu …}.
\end{grammarlens}

\subsection*{Your turn}
\begin{itemize}
\raggedright
  \item \emph{Say it:} \emph{shǒuxiān}
  \item \emph{Say it:} \emph{shǒuxiān}, once more
  \item \emph{Say it:} \emph{shǒuxiān wǒ xiǎng shuō xièxie} — First of all, I'd like to say thank you.
  \item \emph{Say it:} \emph{shǒuxiān wǒmen qù mǎi dōngxi} — To begin with, we go shopping.
  \item \emph{Recall:} say \zh{先}, then \emph{shǒuxiān}: inside the clause, then at its front
\end{itemize}

\begin{tcolorbox}[breakable,colback=teal!4,colframe=teal!35!black,title={Before you move on}]
What does \emph{shǒuxiān} mean? (\textquotedblleft{}First of all\textquotedblright{}.) Where does it stand? (\textbf{At the very front, before the person}.)
\end{tcolorbox}
\section[jiēzhe]{jiēzhe — next; straight after that}

\label{lesson:ZH-C278-jiezhe}



\begin{quote}
Before the new one: last lesson's word, and the pattern it opens.

\begin{itemize}
\raggedright
  \item \emph{Recall:} say \emph{shǒuxiān}, first of all — \textbf{R1}, one lesson back
  \item \emph{Read it:} \zh{先} … \zh{然后} … \zh{最后}, aloud — \emph{xiān … ránhòu … zuìhòu}
\end{itemize}
\end{quote}

\subsection*{You'll want to know: jiēzhe}
\emph{jiēzhe} — \textquotedblleft{}next; straight after that\textquotedblright{}.

\begin{quote}
\emph{wǒ xiān xǐ yīfu, jiēzhe zuòfàn} — First I wash the clothes, and straight after that I cook.
\end{quote}

Its characters come later; for now it is a word for the ear.

\begin{cousinweb}
\emph{jiēzhe} is built from jiē 'to join on' and a light ending, zhe. The first is inherited from Middle Chinese \emph{tsjep} (Guangyun, in Baxter's transcription), glossed 'to join' and 'to hold'. The second is written with a character the dictionary lists under five readings, \emph{drjo}, \emph{trjoX}, \emph{trjoH}, \emph{trjak} and \emph{drjak}, and it glosses none of them as this ending, so the history of that part is not settled.
\end{cousinweb}

\begin{grammarlens}[title={What you've built}]
\emph{ránhòu} leaves a little room between the steps; \emph{jiēzhe} says the next one follows on at once.

Like \zh{然后}, it opens its clause: \emph{wǒ xiān xǐ yīfu, jiēzhe zuòfàn}.
\end{grammarlens}

\subsection*{Your turn}
\begin{itemize}
\raggedright
  \item \emph{Say it:} \emph{jiēzhe}
  \item \emph{Say it:} \emph{jiēzhe}, once more
  \item \emph{Say it:} \emph{wǒ xiān xǐ yīfu, jiēzhe zuòfàn} — First I wash the clothes, and straight after that I cook.
  \item \emph{Say it:} \emph{wǒ xiān chī zǎofàn, jiēzhe qù xuéxiào} — I have breakfast first, then go straight to school.
  \item \emph{Recall:} say \emph{ránhòu}, then \emph{jiēzhe}: then, and straight after
\end{itemize}

\begin{tcolorbox}[breakable,colback=teal!4,colframe=teal!35!black,title={Before you move on}]
What does \emph{jiēzhe} mean? (\textquotedblleft{}Next; straight after that\textquotedblright{}.) Where does it go? (\textbf{At the front of its clause, as \zh{然后} does}.)
\end{tcolorbox}
\section[\texorpdfstring{\zh{吃了饭以后} (chī le fàn yǐhòu)}{吃了饭以后 (chī le fàn yǐhòu)}]{\zh{吃了饭以后} — after doing something}

\label{lesson:ZH-C278-le-yihou}



\begin{quote}
Two things you have, put together in a new way.

\begin{itemize}
\raggedright
  \item \emph{Recall:} say \emph{jiēzhe}, straight after — \textbf{R1}, one lesson back
  \item \emph{Recall:} say \emph{shǒuxiān}, first of all — \textbf{R2}
  \item \emph{Recall:} say \zh{以后} — \emph{yǐhòu}, after; later — \textbf{R3}, durable
  \item \emph{Recall:} say \zh{我看了书} — \emph{wǒ kàn le shū} — \textbf{R3}
\end{itemize}
\end{quote}

\begin{grammarlens}[title={Grammar Lens: \zh{了} … \zh{以后}, after doing something}]
\begin{quote}
\textbf{To say one thing comes after another is done, put \zh{了} after the first verb and \zh{以后} at the end of its clause; the next thing follows.}
\end{quote}

\begin{quote}
\textbf{\zh{吃了饭以后} \zh{我们去买东西}} \emph{chī le fàn yǐhòu, wǒmen qù mǎi dōngxi} — After we've eaten, we'll go shopping. \textbf{\zh{下了课以后} \zh{他去打工}} \emph{xià le kè yǐhòu, tā qù dǎgōng} — After class he goes to his part-time job.
\end{quote}

\zh{了} here is the \zh{了} of a finished action, the one in \zh{我看了书}: the first event is complete before the second starts. It works for what happened and for what is planned; the rest of the sentence tells you which.

\zh{下课} opens up to take \zh{了} in the middle: \zh{下了课}.
\end{grammarlens}

\subsection*{Your turn}
Say these aloud:

\begin{itemize}
\raggedright
  \item \zh{吃了饭以后} \zh{我们去买东西} — \emph{chī le fàn yǐhòu, wǒmen qù mǎi dōngxi} — After we've eaten, we'll go shopping.
  \item \zh{下了课以后} \zh{他去打工} — \emph{xià le kè yǐhòu, tā qù dǎgōng} — After class he goes to his part-time job.
  \item \zh{回了家以后} \zh{我想看书} — \emph{huí le jiā yǐhòu, wǒ xiǎng kàn shū} — Once I'm home, I want to read.
  \item what you do after you have eaten tonight
\end{itemize}

\begin{tcolorbox}[breakable,colback=teal!4,colframe=teal!35!black,title={Before you move on}]
Where does \zh{了} go? (\textbf{After the first verb}.) And \zh{以后}? (\textbf{At the end of the first clause}.) Is the first thing done before the second starts? (\textbf{Yes: that is what \zh{了} says}.)
\end{tcolorbox}
\section[Practice]{\zh{先洗手} \zh{然后吃早饭}}

\label{lesson:ZH-C278-practice}



\begin{quote}
The chapter closes here, with one person's day told in order.

\begin{itemize}
\raggedright
  \item \emph{Read it:} \zh{吃了饭以后}, aloud — \emph{chī le fàn yǐhòu}
  \item \emph{Recall:} say \emph{jiēzhe}, straight after — \textbf{R2}
  \item \emph{Recall:} say \emph{shǒuxiān}, first of all — \textbf{R2}
\end{itemize}
\end{quote}

\subsection*{The exchange}
\begin{quote}
\textbf{A:} \zh{你早上做什么}
\end{quote}

\begin{quote}
\emph{nǐ zǎoshang zuò shénme} — What do you do in the morning?
\end{quote}

\begin{quote}
\textbf{B:} \zh{我先洗手} \zh{然后吃早饭}
\end{quote}

\begin{quote}
\emph{wǒ xiān xǐ shǒu, ránhòu chī zǎofàn} — I wash my hands first, then have breakfast.
\end{quote}

\begin{quote}
\textbf{A:} \zh{吃了早饭以后呢}
\end{quote}

\begin{quote}
\emph{chī le zǎofàn yǐhòu ne} — And after breakfast?
\end{quote}

\begin{quote}
\textbf{B:} \emph{jiēzhe qù xuéxiào} — Straight off to school.
\end{quote}

\begin{quote}
\textbf{A:} \zh{下了课以后呢}
\end{quote}

\begin{quote}
\emph{xià le kè yǐhòu ne} — And after class?
\end{quote}

\begin{quote}
\textbf{B:} \zh{先去打工} \zh{最后回家}
\end{quote}

\begin{quote}
\emph{xiān qù dǎgōng, zuìhòu huí jiā} — First my part-time job, and last home.
\end{quote}

The order of the clauses is the order of the day. B never says \zh{我} after the first line: the person is clear, so Mandarin leaves it out.

\begin{tcolorbox}[breakable,skin=enhanced,colback=orange!6,colframe=orange!45!black,boxrule=0.5pt,arc=1mm,left=6pt,right=6pt,top=4pt,bottom=4pt,fonttitle=\bfseries,title={Read this}]
\emph{Read it:} the exchange through once without stopping, then answer the question below
\end{tcolorbox}

\begin{grammarlens}[title={What you've built — telling it in order}]
Five ways to keep a story in order, all at the front of a clause or right before the verb:

\begin{itemize}
\raggedright
  \item \textbf{\zh{先}} \emph{xiān} — first, before the verb.
  \item \emph{shǒuxiān} — first of all, at the very front.
  \item \textbf{\zh{然后}} \emph{ránhòu} and \emph{jiēzhe} — then, and straight after.
  \item \textbf{\zh{最后}} \emph{zuìhòu} — last.
  \item \textbf{\zh{吃了饭以后}} \emph{chī le fàn yǐhòu} — after doing something.
\end{itemize}

\textbf{Listen.} \emph{Nǐ zuótiān xiān zuò le shénme?}

\textbf{Speak.} Tell your morning in four steps: \zh{先}, \zh{然后}, \emph{jiēzhe}, \zh{最后}.

\emph{Write it:} \zh{最后} — \zh{日} first, then the six strokes, then the piece from the left of \zh{对}
\end{grammarlens}

\subsection*{Your turn}
\begin{itemize}
\raggedright
  \item \emph{Say it:} \zh{我先洗手} \zh{然后吃早饭} \zh{最后去学校} — \emph{wǒ xiān xǐ shǒu, ránhòu chī zǎofàn, zuìhòu qù xuéxiào} — I wash my hands, then have breakfast, and last go to school.
  \item \emph{Say it:} \zh{吃了早饭以后} \zh{我去学校} — \emph{chī le zǎofàn yǐhòu, wǒ qù xuéxiào} — After breakfast I go to school.
  \item \emph{Say it:} the whole exchange, both parts
  \item \emph{Write it:} \zh{先吃饭} \zh{然后看书}
\end{itemize}

\begin{tcolorbox}[breakable,colback=teal!4,colframe=teal!35!black,title={Before you move on}]
Say you wash your hands first, then eat. (\textbf{\zh{我先洗手} \zh{然后吃饭}}.) Say after class. (\textbf{\zh{下了课以后}}.) Which word opens the last step? (\textbf{\zh{最后}}.)
\end{tcolorbox}
